%% kapa: exam=1 round=22 year=2011 period=1 form=A subject=civil_law subject_ko=민법 seq=14 number=14
%% source: Q-Net 감정평가사 기출문제 (cst00309) artlSeq=5036790 file=extracted/1008/5036790/2011년 제22회 감정평가사 1차시험 기출문제/1교시.hwp
%% conversion: kordoc 4.24.51 -> tools/gen1.py (auto); review=proofread-sonnet; flags=none
\begin{question}{14}
집합건물에 관한 설명으로 옳지 않은 것은? (다툼이 있으면 판례에 의함)
\begin{choices}
  \item 구분소유자의 특별승계인에게 그의 체납관리비를 승계하도록 한 관리규약은 공용부분 관리비에 한하여 유효하다.
  \item 관리규약에 구분소유자의 특별승계인에게 그의 체납관리비를 승계하도록 되어 있더라도 관리비에 대한 연체료는 특별승계인에게 승계되지 않는다.
  \item 공용부분의 변경에 관한 사항은 원칙적으로 관리단집회에서 구분소유자의 4분의 3 이상 및 의결권의 4분의 3 이상의 결의로써 결정한다.
  \item 구분소유자는 원칙적으로 그가 가지는 전유부분과 분리하여 대지사용권을 처분할 수 있다.
  \item 공용부분에 대한 공유자의 지분은 그가 가지는 전유부분의 처분에 따른다.
\end{choices}
\end{question}
\answer{4}
